\chapter{Gaussian Mixture Models}

\section{Ziel dieses Kapitels}

Im vorherigen Kapitel haben wir Clustering betrachtet.
Besonders wichtig war \(k\)-means.
\(k\)-means ordnet jeden Datenpunkt genau einem Cluster zu.
Diese Zuordnung ist hart:

\[
    r_{nk}\in\{0,1\}.
\]

Ein Datenpunkt gehört also vollständig zu genau einem Cluster.

Gaussian Mixture Models, kurz GMMs, verallgemeinern diese Idee.
Sie modellieren Daten als Mischung mehrerer Gauß-Verteilungen.
Ein Datenpunkt kann dabei mit unterschiedlichen Wahrscheinlichkeiten zu mehreren Komponenten gehören.

\begin{conceptbox}
Ein Gaussian Mixture Model (GMM) ist ein probabilistisches Modell, das eine Datenverteilung als gewichtete Summe mehrerer Gauß-Verteilungen beschreibt:

\[
    p(\vec{x})
    =
    \sum_{k=1}^{K}
    \pi_k
    \mathcal{N}(\vec{x}\mid\vec{\mu}_k,\Sigma_k).
\]

Dabei beschreibt jede Komponente \(k\) einen Cluster.
\end{conceptbox}

Nach diesem Kapitel solltest du folgende Fragen beantworten können:

\begin{itemize}
    \item Was ist ein Gaussian Mixture Model?
    \item Was ist der Unterschied zwischen \(k\)-means und GMMs?
    \item Was sind Mischgewichte (mixing coefficients)?
    \item Was sind latente Variablen (latent variables)?
    \item Wie schreibt man die Likelihood eines GMMs?
    \item Warum ist die direkte Maximum-Likelihood-Schätzung schwierig?
    \item Was sind Verantwortlichkeiten (responsibilities)?
    \item Wie berechnet man \(p(z=k\mid\vec{x})\)?
    \item Wie funktioniert der Expectation-Maximization-Algorithmus?
    \item Was passiert im E-Schritt?
    \item Was passiert im M-Schritt?
    \item Wie werden \(\vec{\mu}_k\), \(\Sigma_k\) und \(\pi_k\) aktualisiert?
    \item Wie hängt \(k\)-means mit GMMs zusammen?
    \item Welche Probleme können beim Fitten von GMMs auftreten?
\end{itemize}

\section{Motivation}

Viele reale Datensätze bestehen nicht aus einer einzigen einfachen Verteilung.
Stattdessen können verschiedene Gruppen oder Unterpopulationen vorhanden sein.

Beispiele:

\begin{itemize}
    \item Körpergrößen verschiedener Altersgruppen,
    \item medizinische Messwerte verschiedener Patientengruppen,
    \item Bildfeatures verschiedener Objektarten,
    \item Sensordaten verschiedener Betriebszustände,
    \item Kundenverhalten verschiedener Kundensegmente.
\end{itemize}

Eine einzelne Gauß-Verteilung kann solche Daten oft nicht gut beschreiben.
Sie hat nur einen Mittelwert und eine Kovarianzmatrix.

\[
    p(\vec{x})
    =
    \mathcal{N}(\vec{x}\mid\vec{\mu},\Sigma).
\]

Wenn die Daten mehrere Gruppen enthalten, ist eine Mischung mehrerer Gauß-Verteilungen sinnvoller:

\[
    p(\vec{x})
    =
    \sum_{k=1}^{K}
    \pi_k
    \mathcal{N}(\vec{x}\mid\vec{\mu}_k,\Sigma_k).
\]

\begin{conceptbox}
Ein GMM kann multimodale Daten modellieren.

Multimodal bedeutet:
Die Verteilung kann mehrere Häufungsbereiche oder Gipfel haben.
\end{conceptbox}

\section{Eine einzelne multivariate Gauß-Verteilung}

Bevor wir Mischungen betrachten, erinnern wir an die multivariate Gauß-Verteilung.

Für

\[
    \vec{x}\in\mathbb{R}^D
\]

ist die Dichte:

\[
    \mathcal{N}(\vec{x}\mid\vec{\mu},\Sigma)
    =
    \frac{1}{(2\pi)^{D/2}|\Sigma|^{1/2}}
    \exp
    \left(
        -\frac{1}{2}
        (\vec{x}-\vec{\mu})^{\top}
        \Sigma^{-1}
        (\vec{x}-\vec{\mu})
    \right).
\]

Dabei ist:

\begin{itemize}
    \item \(\vec{\mu}\in\mathbb{R}^D\) der Mittelwert,
    \item \(\Sigma\in\mathbb{R}^{D\times D}\) die Kovarianzmatrix,
    \item \(|\Sigma|\) die Determinante der Kovarianzmatrix.
\end{itemize}

\begin{definitionbox}
Die multivariate Gauß-Verteilung ist durch Mittelwert und Kovarianzmatrix bestimmt:

\[
    \vec{\mu}
    =
    \mathbb{E}[\vec{X}],
\]

\[
    \Sigma
    =
    \mathbb{E}
    \left[
        (\vec{X}-\vec{\mu})(\vec{X}-\vec{\mu})^{\top}
    \right].
\]
\end{definitionbox}

\section{MLE für eine einzelne Gauß-Verteilung}

Gegeben seien Datenpunkte

\[
    \vec{x}_1,\dots,\vec{x}_N
    \in
    \mathbb{R}^D.
\]

Wir nehmen an:

\[
    \vec{x}_n
    \overset{\mathrm{i.i.d.}}{\sim}
    \mathcal{N}(\vec{x}\mid\vec{\mu},\Sigma).
\]

Die Log-Likelihood ist:

\[
    \log p(X\mid\vec{\mu},\Sigma)
    =
    \sum_{n=1}^{N}
    \log
    \mathcal{N}(\vec{x}_n\mid\vec{\mu},\Sigma).
\]

Maximum-Likelihood liefert:

\[
    \hat{\vec{\mu}}
    =
    \frac{1}{N}
    \sum_{n=1}^{N}
    \vec{x}_n,
\]

und

\[
    \hat{\Sigma}
    =
    \frac{1}{N}
    \sum_{n=1}^{N}
    (\vec{x}_n-\hat{\vec{\mu}})
    (\vec{x}_n-\hat{\vec{\mu}})^{\top}.
\]

\begin{definitionbox}
Für eine einzelne Gauß-Verteilung sind die Maximum-Likelihood-Schätzer:

\[
    \hat{\vec{\mu}}
    =
    \frac{1}{N}
    \sum_{n=1}^{N}
    \vec{x}_n,
\]

\[
    \hat{\Sigma}
    =
    \frac{1}{N}
    \sum_{n=1}^{N}
    (\vec{x}_n-\hat{\vec{\mu}})
    (\vec{x}_n-\hat{\vec{\mu}})^{\top}.
\]
\end{definitionbox}

\begin{notebox}
Wie bei der eindimensionalen Normalverteilung verwendet der ML-Schätzer der Kovarianz den Faktor \(\frac{1}{N}\), nicht \(\frac{1}{N-1}\).
\end{notebox}

\section{Gaussian Mixture Model}

Ein Gaussian Mixture Model besteht aus \(K\) Gauß-Komponenten.
Jede Komponente hat eigene Parameter:

\[
    \vec{\mu}_k,
    \qquad
    \Sigma_k,
    \qquad
    \pi_k.
\]

Die Dichte ist:

\[
    p(\vec{x})
    =
    \sum_{k=1}^{K}
    \pi_k
    \mathcal{N}(\vec{x}\mid\vec{\mu}_k,\Sigma_k).
\]

\begin{definitionbox}
Ein Gaussian Mixture Model mit \(K\) Komponenten ist:

\[
    p(\vec{x}\mid\Theta)
    =
    \sum_{k=1}^{K}
    \pi_k
    \mathcal{N}(\vec{x}\mid\vec{\mu}_k,\Sigma_k),
\]

wobei

\[
    \Theta
    =
    \{\pi_k,\vec{\mu}_k,\Sigma_k\}_{k=1}^{K}
\]

alle Modellparameter bezeichnet.
\end{definitionbox}

Die Mischgewichte müssen gültige Wahrscheinlichkeiten sein:

\[
    \pi_k\geq 0
\]

und

\[
    \sum_{k=1}^{K}
    \pi_k=1.
\]

\begin{conceptbox}
Das Mischgewicht \(\pi_k\) beschreibt, wie häufig Komponente \(k\) vor der Beobachtung eines Datenpunkts ausgewählt wird.

Großes \(\pi_k\) bedeutet:
Komponente \(k\) erklärt einen großen Anteil der Daten.
\end{conceptbox}

\section{Generative Interpretation}

Ein GMM kann als generatives Modell verstanden werden.
Ein Datenpunkt entsteht in zwei Schritten:

\begin{enumerate}
    \item Wähle eine Komponente \(z\) aus einer kategorialen Verteilung:
    \[
        p(z=k)=\pi_k.
    \]

    \item Ziehe den Datenpunkt aus der entsprechenden Gauß-Verteilung:
    \[
        p(\vec{x}\mid z=k)
        =
        \mathcal{N}(\vec{x}\mid\vec{\mu}_k,\Sigma_k).
    \]
\end{enumerate}

\begin{definitionbox}
Die latente Variable \(z\) gibt an, aus welcher Komponente ein Datenpunkt erzeugt wurde.

Sie ist latent, weil wir \(z\) in den beobachteten Daten normalerweise nicht sehen.
\end{definitionbox}

Die gemeinsame Verteilung von \(\vec{x}\) und \(z\) ist:

\[
    p(\vec{x},z=k)
    =
    p(z=k)p(\vec{x}\mid z=k).
\]

Also:

\[
    p(\vec{x},z=k)
    =
    \pi_k
    \mathcal{N}(\vec{x}\mid\vec{\mu}_k,\Sigma_k).
\]

Wenn wir \(z\) nicht beobachten, müssen wir über alle möglichen Komponenten summieren:

\[
    p(\vec{x})
    =
    \sum_{k=1}^{K}
    p(\vec{x},z=k).
\]

Damit:

\[
    p(\vec{x})
    =
    \sum_{k=1}^{K}
    \pi_k
    \mathcal{N}(\vec{x}\mid\vec{\mu}_k,\Sigma_k).
\]

\begin{conceptbox}
Die Mischung entsteht durch Marginalisierung der latenten Clusterzugehörigkeit:

\[
    p(\vec{x})
    =
    \sum_{z}
    p(\vec{x},z).
\]
\end{conceptbox}

\section{One-Hot-Darstellung der latenten Variablen}

Man kann die latente Variable auch als One-Hot-Vektor schreiben:

\[
    \vec{z}_n
    =
    (z_{n1},\dots,z_{nK})^{\top}.
\]

Dabei gilt:

\[
    z_{nk}\in\{0,1\},
\]

und

\[
    \sum_{k=1}^{K}
    z_{nk}=1.
\]

Wenn

\[
    z_{nk}=1,
\]

dann wurde Datenpunkt \(n\) von Komponente \(k\) erzeugt.

Die Prior-Verteilung über \(\vec{z}_n\) ist:

\[
    p(\vec{z}_n)
    =
    \prod_{k=1}^{K}
    \pi_k^{z_{nk}}.
\]

Die bedingte Verteilung des Datenpunkts ist:

\[
    p(\vec{x}_n\mid\vec{z}_n)
    =
    \prod_{k=1}^{K}
    \mathcal{N}(\vec{x}_n\mid\vec{\mu}_k,\Sigma_k)^{z_{nk}}.
\]

\begin{notebox}
Die One-Hot-Schreibweise ist praktisch für Herleitungen des EM-Algorithmus.

Da genau ein \(z_{nk}\) gleich \(1\) ist, wird genau eine Komponente ausgewählt.
\end{notebox}

\section{Likelihood eines GMMs}

Für i.i.d. Daten

\[
    X
    =
    \{\vec{x}_1,\dots,\vec{x}_N\}
\]

ist die Likelihood:

\[
    p(X\mid\Theta)
    =
    \prod_{n=1}^{N}
    p(\vec{x}_n\mid\Theta).
\]

Da

\[
    p(\vec{x}_n\mid\Theta)
    =
    \sum_{k=1}^{K}
    \pi_k
    \mathcal{N}(\vec{x}_n\mid\vec{\mu}_k,\Sigma_k),
\]

ist:

\[
    p(X\mid\Theta)
    =
    \prod_{n=1}^{N}
    \sum_{k=1}^{K}
    \pi_k
    \mathcal{N}(\vec{x}_n\mid\vec{\mu}_k,\Sigma_k).
\]

\begin{definitionbox}
Die GMM-Likelihood ist:

\[
    p(X\mid\Theta)
    =
    \prod_{n=1}^{N}
    \sum_{k=1}^{K}
    \pi_k
    \mathcal{N}(\vec{x}_n\mid\vec{\mu}_k,\Sigma_k).
\]
\end{definitionbox}

Die Log-Likelihood ist:

\[
    \log p(X\mid\Theta)
    =
    \sum_{n=1}^{N}
    \log
    \left(
        \sum_{k=1}^{K}
        \pi_k
        \mathcal{N}(\vec{x}_n\mid\vec{\mu}_k,\Sigma_k)
    \right).
\]

\begin{definitionbox}
Die GMM-Log-Likelihood ist:

\[
    \ell(\Theta)
    =
    \sum_{n=1}^{N}
    \log
    \left(
        \sum_{k=1}^{K}
        \pi_k
        \mathcal{N}(\vec{x}_n\mid\vec{\mu}_k,\Sigma_k)
    \right).
\]
\end{definitionbox}

\section{Warum direkte MLE schwierig ist}

Bei einer einzelnen Gauß-Verteilung konnten wir die Log-Likelihood direkt ableiten und geschlossene Lösungen erhalten.

Beim GMM steht aber ein Logarithmus über einer Summe:

\[
    \log
    \left(
        \sum_{k=1}^{K}
        \pi_k
        \mathcal{N}(\vec{x}_n\mid\vec{\mu}_k,\Sigma_k)
    \right).
\]

Dieser Ausdruck koppelt die Komponenten.
Dadurch entstehen keine einfachen geschlossenen Lösungen durch direktes Ableiten.

\begin{conceptbox}
Das Problem bei GMMs ist:

\[
    \log
    \left(
        \sum_k
        \text{Komponente}_k
    \right).
\]

Die Summe innerhalb des Logarithmus macht direkte Maximum-Likelihood-Schätzung schwierig.
\end{conceptbox}

Wenn wir die latenten Zuordnungen \(z_n\) kennen würden, wäre das Problem viel einfacher.
Dann könnten wir für jede Komponente separat Mittelwert, Kovarianz und Mischgewicht schätzen.

Da \(z_n\) aber unbekannt ist, verwenden wir Expectation-Maximization.

\section{Posterior über Komponenten}

Für einen Datenpunkt \(\vec{x}_n\) wollen wir wissen:
Wie wahrscheinlich ist es, dass Komponente \(k\) diesen Punkt erzeugt hat?

Das ist:

\[
    p(z_n=k\mid\vec{x}_n).
\]

Mit Bayes' Regel:

\[
    p(z_n=k\mid\vec{x}_n)
    =
    \frac{
        p(z_n=k)p(\vec{x}_n\mid z_n=k)
    }{
        p(\vec{x}_n)
    }.
\]

Einsetzen:

\[
    p(z_n=k\mid\vec{x}_n)
    =
    \frac{
        \pi_k
        \mathcal{N}(\vec{x}_n\mid\vec{\mu}_k,\Sigma_k)
    }{
        \sum_{j=1}^{K}
        \pi_j
        \mathcal{N}(\vec{x}_n\mid\vec{\mu}_j,\Sigma_j)
    }.
\]

Diese Größe heißt Verantwortung oder Responsibility.

\begin{definitionbox}
Die Responsibility \(\gamma_{nk}\) ist die Posterior-Wahrscheinlichkeit, dass Komponente \(k\) für Datenpunkt \(n\) verantwortlich ist:

\[
    \gamma_{nk}
    =
    p(z_n=k\mid\vec{x}_n,\Theta)
    =
    \frac{
        \pi_k
        \mathcal{N}(\vec{x}_n\mid\vec{\mu}_k,\Sigma_k)
    }{
        \sum_{j=1}^{K}
        \pi_j
        \mathcal{N}(\vec{x}_n\mid\vec{\mu}_j,\Sigma_j)
    }.
\]
\end{definitionbox}

Für jeden Datenpunkt gilt:

\[
    \sum_{k=1}^{K}
    \gamma_{nk}=1.
\]

Außerdem:

\[
    0\leq\gamma_{nk}\leq 1.
\]

\begin{conceptbox}
Responsibilities sind weiche Clusterzuordnungen.

Statt

\[
    r_{nk}\in\{0,1\}
\]

wie bei \(k\)-means haben wir

\[
    \gamma_{nk}\in[0,1].
\]
\end{conceptbox}

\section{Effektive Clustergröße}

Die effektive Anzahl der Datenpunkte, die Komponente \(k\) zugeordnet sind, ist:

\[
    N_k
    =
    \sum_{n=1}^{N}
    \gamma_{nk}.
\]

\begin{definitionbox}
Die effektive Clustergröße einer GMM-Komponente ist:

\[
    N_k
    =
    \sum_{n=1}^{N}
    \gamma_{nk}.
\]
\end{definitionbox}

Wenn alle Zuordnungen hart wären, wäre \(N_k\) einfach die Anzahl der Punkte im Cluster.
Bei weichen Zuordnungen zählt jeder Punkt anteilig.

\begin{examplebox}
Wenn ein Punkt Responsibility

\[
    \gamma_{n1}=0{,}7
\]

für Komponente \(1\) und

\[
    \gamma_{n2}=0{,}3
\]

für Komponente \(2\) hat, dann zählt er zu \(70\%\) zu Komponente \(1\) und zu \(30\%\) zu Komponente \(2\).
\end{examplebox}

\section{Expectation-Maximization: Grundidee}

Expectation-Maximization, kurz EM, ist ein Algorithmus für Modelle mit latenten Variablen.

Die Grundidee:

\begin{enumerate}
    \item E-Schritt:
    Schätze die unbekannten latenten Zuordnungen mit den aktuellen Parametern.

    \item M-Schritt:
    Aktualisiere die Parameter so, als wären diese geschätzten Zuordnungen bekannt.
\end{enumerate}

\begin{conceptbox}
EM wechselt zwischen:

\[
    \text{Erwarte latente Struktur}
    \quad
    \longrightarrow
    \quad
    \text{Maximiere Parameter}.
\]

Für GMMs bedeutet das:

\[
    \text{E-Schritt: Responsibilities berechnen}.
\]

\[
    \text{M-Schritt: Gauß-Parameter mit Responsibilities neu schätzen}.
\]
\end{conceptbox}

\section{E-Schritt für GMMs}

Im E-Schritt werden die Responsibilities mit den aktuellen Parametern berechnet.
Seien die aktuellen Parameter:

\[
    \Theta^{\mathrm{old}}
    =
    \{\pi_k^{\mathrm{old}},\vec{\mu}_k^{\mathrm{old}},\Sigma_k^{\mathrm{old}}\}_{k=1}^{K}.
\]

Dann ist:

\[
    \gamma_{nk}
    =
    \frac{
        \pi_k^{\mathrm{old}}
        \mathcal{N}(\vec{x}_n\mid\vec{\mu}_k^{\mathrm{old}},\Sigma_k^{\mathrm{old}})
    }{
        \sum_{j=1}^{K}
        \pi_j^{\mathrm{old}}
        \mathcal{N}(\vec{x}_n\mid\vec{\mu}_j^{\mathrm{old}},\Sigma_j^{\mathrm{old}})
    }.
\]

\begin{definitionbox}
Der E-Schritt eines GMMs berechnet für alle Datenpunkte und Komponenten:

\[
    \gamma_{nk}
    =
    p(z_n=k\mid\vec{x}_n,\Theta^{\mathrm{old}}).
\]
\end{definitionbox}

\section{M-Schritt für GMMs}

Im M-Schritt werden die Parameter mit den Responsibilities aktualisiert.

Zuerst:

\[
    N_k
    =
    \sum_{n=1}^{N}
    \gamma_{nk}.
\]

Der neue Mittelwert ist der gewichtete Mittelwert:

\[
    \vec{\mu}_k^{\mathrm{new}}
    =
    \frac{1}{N_k}
    \sum_{n=1}^{N}
    \gamma_{nk}
    \vec{x}_n.
\]

Die neue Kovarianzmatrix ist:

\[
    \Sigma_k^{\mathrm{new}}
    =
    \frac{1}{N_k}
    \sum_{n=1}^{N}
    \gamma_{nk}
    (\vec{x}_n-\vec{\mu}_k^{\mathrm{new}})
    (\vec{x}_n-\vec{\mu}_k^{\mathrm{new}})^{\top}.
\]

Das neue Mischgewicht ist:

\[
    \pi_k^{\mathrm{new}}
    =
    \frac{N_k}{N}.
\]

\begin{definitionbox}
Der M-Schritt eines GMMs aktualisiert:

\[
    N_k
    =
    \sum_{n=1}^{N}
    \gamma_{nk},
\]

\[
    \vec{\mu}_k
    =
    \frac{1}{N_k}
    \sum_{n=1}^{N}
    \gamma_{nk}
    \vec{x}_n,
\]

\[
    \Sigma_k
    =
    \frac{1}{N_k}
    \sum_{n=1}^{N}
    \gamma_{nk}
    (\vec{x}_n-\vec{\mu}_k)
    (\vec{x}_n-\vec{\mu}_k)^{\top},
\]

\[
    \pi_k
    =
    \frac{N_k}{N}.
\]
\end{definitionbox}

\section{Interpretation der M-Schritt-Formeln}

Die Formeln sind die gleichen wie bei einer Gauß-Verteilung, aber mit Gewichten.

Der Mittelwert

\[
    \vec{\mu}_k
\]

ist der gewichtete Durchschnitt aller Datenpunkte.
Ein Punkt mit großer Responsibility \(\gamma_{nk}\) für Komponente \(k\) beeinflusst \(\vec{\mu}_k\) stark.
Ein Punkt mit kleiner Responsibility beeinflusst \(\vec{\mu}_k\) wenig.

Die Kovarianz

\[
    \Sigma_k
\]

ist die gewichtete Kovarianz um den neuen Mittelwert.

Das Mischgewicht

\[
    \pi_k
\]

ist der Anteil der effektiven Punkte, die Komponente \(k\) erklärt.

\begin{conceptbox}
Der M-Schritt ist wie das Fitten einer Gauß-Verteilung pro Cluster, aber jeder Datenpunkt zählt nur anteilig gemäß seiner Responsibility.
\end{conceptbox}

\section{EM-Algorithmus für GMMs}

Der vollständige Algorithmus lautet:

\begin{definitionbox}
EM für Gaussian Mixture Models:

\begin{enumerate}
    \item Wähle \(K\).
    \item Initialisiere \(\pi_k\), \(\vec{\mu}_k\) und \(\Sigma_k\).
    \item E-Schritt:
    \[
        \gamma_{nk}
        =
        \frac{
            \pi_k
            \mathcal{N}(\vec{x}_n\mid\vec{\mu}_k,\Sigma_k)
        }{
            \sum_{j=1}^{K}
            \pi_j
            \mathcal{N}(\vec{x}_n\mid\vec{\mu}_j,\Sigma_j)
        }.
    \]
    \item M-Schritt:
    \[
        N_k
        =
        \sum_{n=1}^{N}
        \gamma_{nk},
    \]
    \[
        \vec{\mu}_k
        =
        \frac{1}{N_k}
        \sum_{n=1}^{N}
        \gamma_{nk}\vec{x}_n,
    \]
    \[
        \Sigma_k
        =
        \frac{1}{N_k}
        \sum_{n=1}^{N}
        \gamma_{nk}
        (\vec{x}_n-\vec{\mu}_k)
        (\vec{x}_n-\vec{\mu}_k)^{\top},
    \]
    \[
        \pi_k
        =
        \frac{N_k}{N}.
    \]
    \item Berechne die Log-Likelihood:
    \[
        \ell(\Theta)
        =
        \sum_{n=1}^{N}
        \log
        \left(
            \sum_{k=1}^{K}
            \pi_k
            \mathcal{N}(\vec{x}_n\mid\vec{\mu}_k,\Sigma_k)
        \right).
    \]
    \item Wiederhole E- und M-Schritt, bis die Log-Likelihood kaum noch steigt.
\end{enumerate}
\end{definitionbox}

\section{Konvergenz von EM}

EM hat eine wichtige Eigenschaft:
Die Log-Likelihood sinkt nicht nach einer Iteration.

\[
    \ell(\Theta^{\mathrm{new}})
    \geq
    \ell(\Theta^{\mathrm{old}}).
\]

\begin{conceptbox}
EM verbessert oder erhält in jeder Iteration die Log-Likelihood.

Das bedeutet:
Der Algorithmus bewegt sich schrittweise zu einem lokalen Maximum der Likelihood.
\end{conceptbox}

Wie bei \(k\)-means bedeutet das aber nicht, dass das globale Maximum gefunden wird.
EM kann in lokalen Maxima landen.

\begin{notebox}
Konvergenz von EM bedeutet meistens:
Die Änderung der Log-Likelihood wird kleiner als eine gewählte Toleranz.

Zum Beispiel:

\[
    |\ell^{(\tau+1)}-\ell^{(\tau)}|<\varepsilon.
\]
\end{notebox}

\section{Free Energy und EM-Idee}

Eine tiefere Sicht auf EM verwendet eine untere Schranke der Log-Likelihood.
Diese Schranke wird oft Free Energy genannt.

Für GMMs kann man schreiben:

\[
    \mathcal{F}
    =
    \sum_{n=1}^{N}
    \sum_{k=1}^{K}
    \gamma_{nk}
    \left(
        \log \pi_k
        +
        \log
        \mathcal{N}(\vec{x}_n\mid\vec{\mu}_k,\Sigma_k)
    \right)
    -
    \sum_{n=1}^{N}
    \sum_{k=1}^{K}
    \gamma_{nk}
    \log \gamma_{nk}.
\]

\begin{definitionbox}
Die Free Energy für ein GMM ist:

\[
    \mathcal{F}
    =
    \sum_{n=1}^{N}
    \sum_{k=1}^{K}
    \gamma_{nk}
    \left(
        \log \pi_k
        +
        \log
        \mathcal{N}(\vec{x}_n\mid\vec{\mu}_k,\Sigma_k)
    \right)
    -
    \sum_{n=1}^{N}
    \sum_{k=1}^{K}
    \gamma_{nk}
    \log \gamma_{nk}.
\]
\end{definitionbox}

Die E- und M-Schritte können so interpretiert werden:

\begin{itemize}
    \item E-Schritt:
    Wähle \(\gamma_{nk}\), sodass die Schranke eng wird.
    \item M-Schritt:
    Maximiere die Schranke bezüglich der Modellparameter.
\end{itemize}

\begin{notebox}
Für praktische Rechnungen reichen meistens die E- und M-Formeln.
Die Free-Energy-Sicht erklärt aber, warum EM die Log-Likelihood nicht verringert.
\end{notebox}

\section{Initialisierung}

GMMs sind empfindlich gegenüber Initialisierung.
Man muss initialisieren:

\[
    \vec{\mu}_k,
    \qquad
    \Sigma_k,
    \qquad
    \pi_k.
\]

Typische Strategien:

\begin{itemize}
    \item Mittelwerte zufällig aus Datenpunkten wählen,
    \item \(k\)-means verwenden und die Zentren übernehmen,
    \item Kovarianzmatrizen initial als empirische Kovarianz setzen,
    \item Mischgewichte gleich verteilen:
    \[
        \pi_k=\frac{1}{K}.
    \]
\end{itemize}

\begin{conceptbox}
Eine häufige Praxis ist:
Erst \(k\)-means ausführen, dann die Zentren und Zuordnungen zur Initialisierung des GMMs verwenden.
\end{conceptbox}

\section{Vergleich mit \(k\)-means}

\(k\)-means und GMMs sind eng verwandt.

\[
\begin{array}{c|c|c}
\text{Eigenschaft}
&
k\text{-means}
&
\text{GMM}
\\
\hline
\text{Zuordnung}
&
\text{hart}
&
\text{weich}
\\
\hline
\text{Clusterform}
&
\text{kugelförmig}
&
\text{elliptisch möglich}
\\
\hline
\text{Parameter}
&
\vec{\mu}_k
&
\pi_k,\vec{\mu}_k,\Sigma_k
\\
\hline
\text{Optimierung}
&
\text{alternierend}
&
\text{EM}
\\
\hline
\text{Probabilistisch}
&
\text{nicht direkt}
&
\text{ja}
\end{array}
\]

\begin{conceptbox}
GMMs verallgemeinern \(k\)-means:

\begin{itemize}
    \item Cluster haben Kovarianzmatrizen.
    \item Punkte können mehreren Clustern anteilig zugeordnet sein.
    \item Das Modell liefert Wahrscheinlichkeitsdichten.
\end{itemize}
\end{conceptbox}

\section{\(k\)-means als Grenzfall eines GMMs}

Man kann \(k\)-means als Grenzfall eines GMMs verstehen.
Angenommen, alle Komponenten haben gleiche isotrope Kovarianz:

\[
    \Sigma_k=\sigma^2 I.
\]

Wenn

\[
    \sigma^2\to 0,
\]

dann wird die Gauß-Dichte um jedes Zentrum sehr schmal.
Jeder Punkt wird fast vollständig der nächsten Komponente zugeordnet.

Dann werden die Responsibilities fast hart:

\[
    \gamma_{nk}
    \approx
    \begin{cases}
        1, & k=\arg\min_j\|\vec{x}_n-\vec{\mu}_j\|^2, \\
        0, & \text{sonst}.
    \end{cases}
\]

\begin{conceptbox}
\(k\)-means entspricht ungefähr einem GMM mit gleichen kugelförmigen Kovarianzen und harten Zuordnungen.
\end{conceptbox}

\section{Kovarianztypen}

In GMMs kann man verschiedene Formen für \(\Sigma_k\) verwenden.

\subsection{Volle Kovarianz}

Jede Komponente hat eine vollständige Kovarianzmatrix:

\[
    \Sigma_k\in\mathbb{R}^{D\times D}.
\]

Das erlaubt beliebig orientierte elliptische Cluster.

\subsection{Diagonale Kovarianz}

Jede Kovarianzmatrix ist diagonal:

\[
    \Sigma_k
    =
    \mathrm{diag}
    (\sigma_{k1}^2,\dots,\sigma_{kD}^2).
\]

Dann sind die Merkmale innerhalb einer Komponente unkorreliert.

\subsection{Sphärische Kovarianz}

Jede Komponente hat:

\[
    \Sigma_k
    =
    \sigma_k^2 I.
\]

Dann sind Cluster kugelförmig, aber jede Komponente kann eigene Varianz haben.

\subsection{Geteilte Kovarianz}

Alle Komponenten teilen dieselbe Kovarianzmatrix:

\[
    \Sigma_k=\Sigma
    \quad
    \text{für alle}
    \quad
    k.
\]

\begin{notebox}
Volle Kovarianzen sind flexibel, benötigen aber viele Parameter.
Diagonale oder sphärische Kovarianzen sind restriktiver, aber stabiler bei wenig Daten.
\end{notebox}

\section{Anzahl der Parameter}

Die Anzahl der Parameter eines GMMs wächst mit \(K\) und \(D\).

Für jede Komponente:

\begin{itemize}
    \item Mittelwert:
    \[
        D
    \]
    Parameter.
    \item Volle Kovarianzmatrix:
    \[
        \frac{D(D+1)}{2}
    \]
    Parameter, weil sie symmetrisch ist.
    \item Mischgewicht:
    Ein Gewicht pro Komponente, aber wegen
    \[
        \sum_{k=1}^{K}\pi_k=1
    \]
    gibt es insgesamt nur \(K-1\) freie Mischgewichtsparameter.
\end{itemize}

\begin{conceptbox}
Ein GMM mit voller Kovarianz kann in hoher Dimension sehr viele Parameter haben.

Daher braucht man genügend Daten oder Einschränkungen der Kovarianzstruktur.
\end{conceptbox}

\section{Wahl der Komponentenzahl \(K\)}

Die Anzahl der Komponenten \(K\) ist ein Hyperparameter.
Zu kleines \(K\) führt zu einem zu groben Modell.
Zu großes \(K\) kann Overfitting erzeugen.

Mögliche Methoden zur Wahl von \(K\):

\begin{itemize}
    \item Validierungs-Log-Likelihood,
    \item AIC,
    \item BIC,
    \item fachliche Interpretierbarkeit,
    \item Vergleich mehrerer Initialisierungen.
\end{itemize}

\begin{definitionbox}
Die Bayesian Information Criterion, kurz BIC, ist:

\[
    \mathrm{BIC}
    =
    -2\ell(\hat{\Theta})
    +
    p\log N,
\]

wobei \(p\) die Anzahl freier Parameter und \(N\) die Anzahl der Datenpunkte ist.
\end{definitionbox}

Ein kleinerer BIC-Wert ist besser.
Der erste Term belohnt gute Likelihood.
Der zweite Term bestraft viele Parameter.

\begin{notebox}
BIC ist ein Modellselektionskriterium.
Es ersetzt nicht automatisch fachliche Prüfung, kann aber bei der Wahl von \(K\) helfen.
\end{notebox}

\section{Clustering mit GMMs}

Nach dem Fitten eines GMMs kann man Datenpunkte clustern.

Die weiche Zuordnung ist:

\[
    \gamma_{nk}
    =
    p(z_n=k\mid\vec{x}_n).
\]

Eine harte Zuordnung erhält man durch:

\[
    \hat{z}_n
    =
    \arg\max_k
    \gamma_{nk}.
\]

\begin{definitionbox}
Die harte GMM-Clusterzuordnung ist:

\[
    \hat{z}_n
    =
    \arg\max_k
    \gamma_{nk}.
\]
\end{definitionbox}

\begin{conceptbox}
GMMs liefern mehr Information als \(k\)-means:

\begin{itemize}
    \item Die wahrscheinlichste Komponente.
    \item Die Unsicherheit der Zuordnung.
    \item Eine Dichte \(p(\vec{x})\) für neue Datenpunkte.
\end{itemize}
\end{conceptbox}

\section{Anomalieerkennung mit GMMs}

Da ein GMM eine Dichte \(p(\vec{x})\) modelliert, kann es auch für Anomalieerkennung verwendet werden.

Punkte mit sehr kleiner Dichte sind ungewöhnlich:

\[
    p(\vec{x}) \text{ klein}
    \quad
    \Rightarrow
    \quad
    \vec{x} \text{ ist potenziell anomal}.
\]

Oder mit negativer Log-Likelihood:

\[
    -\log p(\vec{x})
\]

als Anomaliescore.

\begin{definitionbox}
Ein GMM-Anomaliescore kann definiert werden als:

\[
    s(\vec{x})
    =
    -\log p(\vec{x}).
\]

Große Werte bedeuten geringe Modellwahrscheinlichkeit.
\end{definitionbox}

\section{Singularitäten und degenerierte Lösungen}

Ein wichtiges Problem bei GMMs:
Die Likelihood kann unbeschränkt werden.

Wenn eine Komponente genau auf einen einzelnen Datenpunkt gesetzt wird und ihre Kovarianz gegen \(0\) geht, wird die Gauß-Dichte an diesem Punkt sehr groß.

Formal:

\[
    \Sigma_k \to 0
\]

kann zu

\[
    \mathcal{N}(\vec{x}_n\mid\vec{\mu}_k,\Sigma_k)\to\infty
\]

führen, wenn

\[
    \vec{\mu}_k=\vec{x}_n.
\]

\begin{examtipbox}
Die GMM-Likelihood kann singuläre Lösungen haben.

Das ist ein wichtiger Unterschied zu vielen einfacheren Modellen.
\end{examtipbox}

Praktische Gegenmaßnahmen:

\begin{itemize}
    \item Kovarianzmatrizen regularisieren:
    \[
        \Sigma_k
        \leftarrow
        \Sigma_k+\varepsilon I.
    \]
    \item Mindestvarianzen erzwingen.
    \item Mehrere Initialisierungen verwenden.
    \item \(K\) nicht zu groß wählen.
    \item Diagonale oder geteilte Kovarianzen verwenden.
\end{itemize}

\section{Numerische Stabilität}

Die GMM-Dichte kann sehr kleine Werte annehmen.
Bei vielen Datenpunkten können Produkte numerisch unterlaufen.

Deshalb arbeitet man mit Log-Likelihoods:

\[
    \ell(\Theta)
    =
    \sum_{n=1}^{N}
    \log
    \left(
        \sum_{k=1}^{K}
        \pi_k
        \mathcal{N}(\vec{x}_n\mid\vec{\mu}_k,\Sigma_k)
    \right).
\]

Auch der Ausdruck

\[
    \log
    \left(
        \sum_k
        \exp(a_k)
    \right)
\]

sollte numerisch stabil berechnet werden.
Dafür verwendet man die Log-Sum-Exp-Technik.

\begin{definitionbox}
Die Log-Sum-Exp-Identität ist:

\[
    \log\sum_k\exp(a_k)
    =
    m
    +
    \log
    \sum_k
    \exp(a_k-m),
\]

wobei

\[
    m=\max_k a_k.
\]
\end{definitionbox}

\begin{notebox}
Durch Subtraktion von \(m\) vermeidet man, dass \(\exp(a_k)\) numerisch überläuft.
\end{notebox}

\section{GMMs in hoher Dimension}

In hoher Dimension können GMMs schwierig zu fitten sein.

Gründe:

\begin{itemize}
    \item Volle Kovarianzmatrizen haben viele Parameter.
    \item Kovarianzmatrizen können singulär werden.
    \item Dichteschätzung wird durch den Fluch der Dimensionalität schwieriger.
    \item Viele Datenpunkte sind nötig, um Kovarianzen stabil zu schätzen.
\end{itemize}

\begin{conceptbox}
In hoher Dimension verwendet man oft:

\begin{itemize}
    \item PCA vor GMM,
    \item diagonale Kovarianzmatrizen,
    \item Regularisierung der Kovarianzen,
    \item weniger Komponenten,
    \item fachlich motivierte Features.
\end{itemize}
\end{conceptbox}

\section{GMMs und PCA}

PCA und GMMs können kombiniert werden.
PCA reduziert zuerst die Dimension:

\[
    \vec{x}
    \in\mathbb{R}^D
    \quad
    \longrightarrow
    \quad
    \vec{z}
    \in\mathbb{R}^M.
\]

Dann wird ein GMM im PCA-Raum gefittet:

\[
    p(\vec{z})
    =
    \sum_{k=1}^{K}
    \pi_k
    \mathcal{N}(\vec{z}\mid\vec{\mu}_k,\Sigma_k).
\]

\begin{notebox}
Wenn PCA als Vorverarbeitung in einem Workflow verwendet wird, muss sie nur auf Trainingsdaten gefittet werden, um Data Leakage zu vermeiden.
\end{notebox}

\section{Vor- und Nachteile von GMMs}

\subsection{Vorteile}

\begin{itemize}
    \item probabilistische Interpretation,
    \item weiche Clusterzuordnungen,
    \item elliptische Cluster möglich,
    \item Dichte \(p(\vec{x})\) verfügbar,
    \item Anomalieerkennung möglich,
    \item gute Verbindung zu EM und latenten Variablen.
\end{itemize}

\subsection{Nachteile}

\begin{itemize}
    \item Anzahl der Komponenten \(K\) muss gewählt werden,
    \item lokale Optima,
    \item empfindlich gegenüber Initialisierung,
    \item mögliche singuläre Kovarianzen,
    \item viele Parameter bei voller Kovarianz,
    \item Gauß-Annahme nicht immer passend.
\end{itemize}

\begin{conceptbox}
GMMs sind flexibler als \(k\)-means, aber auch empfindlicher und parameterreicher.
\end{conceptbox}

\section{Mini-Aufgaben}

\subsection{Aufgabe 1: Mischgewichte prüfen}

Gegeben seien:

\[
    \pi_1=0{,}2,
    \qquad
    \pi_2=0{,}5,
    \qquad
    \pi_3=0{,}3.
\]

Sind das gültige Mischgewichte?

\begin{examplebox}
Mischgewichte müssen nichtnegativ sein und sich zu \(1\) summieren.

Alle Werte sind nichtnegativ.
Außerdem:

\[
    0{,}2+0{,}5+0{,}3=1.
\]

Also sind die Mischgewichte gültig.
\end{examplebox}

\subsection{Aufgabe 2: GMM-Dichte}

Ein eindimensionales GMM habe zwei Komponenten:

\[
    \pi_1=0{,}4,
    \qquad
    \pi_2=0{,}6.
\]

Die Dichten an einem Punkt \(x\) seien:

\[
    \mathcal{N}(x\mid\mu_1,\sigma_1^2)=0{,}1,
\]

\[
    \mathcal{N}(x\mid\mu_2,\sigma_2^2)=0{,}3.
\]

Berechne \(p(x)\).

\begin{examplebox}
Die GMM-Dichte ist:

\[
    p(x)
    =
    \sum_{k=1}^{2}
    \pi_k
    \mathcal{N}(x\mid\mu_k,\sigma_k^2).
\]

Einsetzen:

\[
    p(x)
    =
    0{,}4\cdot 0{,}1
    +
    0{,}6\cdot 0{,}3.
\]

Also:

\[
    p(x)
    =
    0{,}04+0{,}18
    =
    0{,}22.
\]
\end{examplebox}

\subsection{Aufgabe 3: Responsibility berechnen}

Mit den Werten aus Aufgabe \(2\) berechne die Responsibility von Komponente \(2\).

\begin{examplebox}
Die Responsibility ist:

\[
    \gamma_2
    =
    \frac{
        \pi_2
        \mathcal{N}(x\mid\mu_2,\sigma_2^2)
    }{
        \sum_{j=1}^{2}
        \pi_j
        \mathcal{N}(x\mid\mu_j,\sigma_j^2)
    }.
\]

Der Nenner wurde in Aufgabe \(2\) berechnet:

\[
    p(x)=0{,}22.
\]

Der Zähler ist:

\[
    0{,}6\cdot 0{,}3=0{,}18.
\]

Also:

\[
    \gamma_2
    =
    \frac{0{,}18}{0{,}22}
    =
    \frac{18}{22}
    =
    \frac{9}{11}
    \approx
    0{,}818.
\]
\end{examplebox}

\subsection{Aufgabe 4: Effektive Clustergröße}

Für eine Komponente seien die Responsibilities über vier Datenpunkte:

\[
    \gamma_{1k}=0{,}8,
    \quad
    \gamma_{2k}=0{,}6,
    \quad
    \gamma_{3k}=0{,}1,
    \quad
    \gamma_{4k}=0{,}5.
\]

Berechne \(N_k\).

\begin{examplebox}
Die effektive Clustergröße ist:

\[
    N_k
    =
    \sum_{n=1}^{N}
    \gamma_{nk}.
\]

Also:

\[
    N_k
    =
    0{,}8+0{,}6+0{,}1+0{,}5
    =
    2{,}0.
\]
\end{examplebox}

\subsection{Aufgabe 5: Mischgewicht aktualisieren}

Für ein GMM mit

\[
    N=10
\]

Datenpunkten sei

\[
    N_k=2.
\]

Berechne das neue Mischgewicht \(\pi_k\).

\begin{examplebox}
Im M-Schritt gilt:

\[
    \pi_k
    =
    \frac{N_k}{N}.
\]

Einsetzen:

\[
    \pi_k
    =
    \frac{2}{10}
    =
    0{,}2.
\]
\end{examplebox}

\subsection{Aufgabe 6: GMM oder \(k\)-means?}

Ein Verfahren gibt für einen Datenpunkt folgende Zuordnungen aus:

\[
    \gamma_{n1}=0{,}2,
    \qquad
    \gamma_{n2}=0{,}7,
    \qquad
    \gamma_{n3}=0{,}1.
\]

Handelt es sich eher um \(k\)-means oder um ein GMM?

\begin{examplebox}
Die Werte sind weich und summieren sich zu \(1\):

\[
    0{,}2+0{,}7+0{,}1=1.
\]

Das spricht für ein GMM oder ein anderes Soft-Clustering-Verfahren.

Bei \(k\)-means wäre die Zuordnung hart, also zum Beispiel:

\[
    (0,1,0).
\]
\end{examplebox}

\section{Häufige Fehler}

\begin{examtipbox}
Typische Fehler bei Gaussian Mixture Models:

\begin{enumerate}
    \item GMMs mit \(k\)-means gleichsetzen.
    \item Harte und weiche Zuordnungen verwechseln.
    \item Mischgewichte \(\pi_k\) nicht auf Summe \(1\) normieren.
    \item Responsibilities nicht über \(k\) normieren.
    \item Die GMM-Likelihood ohne die Summe über Komponenten schreiben.
    \item Die Log-Likelihood fälschlich als Summe der Log-Komponenten schreiben.
    \item Im M-Schritt \(N_k=\sum_n\gamma_{nk}\) vergessen.
    \item Kovarianzmatrizen ohne Gewichtung durch \(\gamma_{nk}\) berechnen.
    \item Singuläre Kovarianzen ignorieren.
    \item EM für garantiert global optimal halten.
    \item Zu viele Komponenten ohne Regularisierung verwenden.
    \item In hoher Dimension volle Kovarianzen ohne genügend Daten verwenden.
\end{enumerate}
\end{examtipbox}

\section{Zusammenfassung}

\begin{itemize}
    \item Ein Gaussian Mixture Model modelliert Daten als Mischung von Gauß-Verteilungen:
    \[
        p(\vec{x})
        =
        \sum_{k=1}^{K}
        \pi_k
        \mathcal{N}(\vec{x}\mid\vec{\mu}_k,\Sigma_k).
    \]
    \item Die Mischgewichte erfüllen:
    \[
        \pi_k\geq 0,
        \qquad
        \sum_{k=1}^{K}\pi_k=1.
    \]
    \item Die latente Variable \(z\) gibt an, aus welcher Komponente ein Datenpunkt stammt.
    \item Die gemeinsame Verteilung ist:
    \[
        p(\vec{x},z=k)
        =
        \pi_k
        \mathcal{N}(\vec{x}\mid\vec{\mu}_k,\Sigma_k).
    \]
    \item Die GMM-Log-Likelihood ist:
    \[
        \ell(\Theta)
        =
        \sum_{n=1}^{N}
        \log
        \left(
            \sum_{k=1}^{K}
            \pi_k
            \mathcal{N}(\vec{x}_n\mid\vec{\mu}_k,\Sigma_k)
        \right).
    \]
    \item Die Responsibility ist:
    \[
        \gamma_{nk}
        =
        \frac{
            \pi_k
            \mathcal{N}(\vec{x}_n\mid\vec{\mu}_k,\Sigma_k)
        }{
            \sum_{j=1}^{K}
            \pi_j
            \mathcal{N}(\vec{x}_n\mid\vec{\mu}_j,\Sigma_j)
        }.
    \]
    \item Die effektive Clustergröße ist:
    \[
        N_k
        =
        \sum_{n=1}^{N}
        \gamma_{nk}.
    \]
    \item Im M-Schritt gilt:
    \[
        \vec{\mu}_k
        =
        \frac{1}{N_k}
        \sum_{n=1}^{N}
        \gamma_{nk}\vec{x}_n,
    \]
    \[
        \Sigma_k
        =
        \frac{1}{N_k}
        \sum_{n=1}^{N}
        \gamma_{nk}
        (\vec{x}_n-\vec{\mu}_k)
        (\vec{x}_n-\vec{\mu}_k)^{\top},
    \]
    \[
        \pi_k
        =
        \frac{N_k}{N}.
    \]
    \item EM wechselt zwischen E-Schritt und M-Schritt.
    \item EM erhöht die Log-Likelihood monoton, findet aber nicht garantiert das globale Maximum.
    \item GMMs liefern weiche Zuordnungen und können elliptische Cluster modellieren.
    \item \(k\)-means kann als Grenzfall eines GMMs mit harten Zuordnungen verstanden werden.
    \item GMMs können Probleme mit lokalen Optima, Initialisierung und singulären Kovarianzen haben.
\end{itemize}

\section{Typische Prüfungsfragen}

\begin{examtipbox}
Mögliche Prüfungsfragen zu diesem Kapitel:

\begin{enumerate}
    \item Was ist ein Gaussian Mixture Model?
    \item Welche Parameter hat ein GMM?
    \item Welche Bedingungen müssen die Mischgewichte erfüllen?
    \item Erkläre die generative Interpretation eines GMMs.
    \item Was ist eine latente Variable?
    \item Schreibe \(p(\vec{x},z=k)\) für ein GMM auf.
    \item Leite \(p(\vec{x})\) durch Marginalisierung von \(z\) her.
    \item Schreibe die GMM-Likelihood und Log-Likelihood auf.
    \item Warum ist direkte Maximum-Likelihood-Schätzung bei GMMs schwierig?
    \item Was ist eine Responsibility?
    \item Leite die Responsibility mit Bayes' Regel her.
    \item Was ist \(N_k\)?
    \item Was passiert im E-Schritt?
    \item Was passiert im M-Schritt?
    \item Leite die Update-Formel für \(\vec{\mu}_k\) her.
    \item Leite die Update-Formel für \(\Sigma_k\) her.
    \item Leite die Update-Formel für \(\pi_k\) her.
    \item Warum steigt die Log-Likelihood bei EM monoton?
    \item Warum findet EM nicht garantiert das globale Optimum?
    \item Wie kann man GMMs initialisieren?
    \item Was ist der Unterschied zwischen GMM und \(k\)-means?
    \item In welchem Sinn ist \(k\)-means ein Grenzfall eines GMMs?
    \item Welche Kovarianztypen kann man bei GMMs verwenden?
    \item Warum können GMMs singuläre Lösungen haben?
    \item Wie kann man \(K\) für ein GMM wählen?
    \item Wie kann ein GMM für Anomalieerkennung verwendet werden?
\end{enumerate}
\end{examtipbox}